\section{System Design}

The source of truth is Lean plus explicit ledgers.  A mathematical ingredient
is never treated as proved merely because it is plausible in prose.  If it is
not a compiled Lean declaration, it must be a \ProofObligation{} or a cited
result with a recorded Lean status.

The workflow has two public modes.  In \Lean{faithfulPaper} mode, the theorem,
assumptions, constants, and proof target are fixed by the source paper.  In
\Lean{exploratoryProof} mode, the system tracks proof routes for active
research drafts, such as RMFLD, without promoting candidate assumptions to
facts.

The project-paper framing follows the automation emphasis of
\url{https://github.com/wanshuiyin/Auto-claude-code-research-in-sleep}, the search-memory perspective of \url{https://github.com/Trinkle23897/learning-beyond-gradients}, the
population-and-archive idea of \url{https://github.com/FeiLiu36/EoH}, the Lean-specific blueprint
control loop of \url{https://github.com/YuanheZ/LeanMarathon}, and the Lean proof-diagnostics
workflow exemplified by \url{https://github.com/math-ai-org/mathcode}.  The acceptance condition is
stricter than search success: every cycle must leave a buildable repository or
an explicit blocked status, and post-blueprint cycles must identify the
dynamic leaf or illness area they changed.

\begin{figure}[t]
\centering
\begin{tikzpicture}[
  box/.style={draw, rounded corners=2pt, align=left, minimum width=0.85\textwidth, inner sep=6pt},
  arrow/.style={-{Latex[length=2mm]}, thick}
]
\node[box, fill=purple!5] (upper) {\textbf{Upper}: mode, objective, phase discipline, memory compression};
\node[box, fill=blue!5, below=0.45cm of upper] (middle) {\textbf{Middle}: source-to-Lean and Lean-to-Markdown/LaTeX conversion, proof DAGs, cited results};
\node[box, fill=orange!8, below=0.45cm of middle] (lower) {\textbf{Lower}: one Lean declaration, proof block, source-index change, or obligation refinement};
\node[box, fill=green!6, below=0.45cm of lower] (reviewer) {\textbf{Reviewer}: build gate, source correspondence, no fake proof closure, phase discipline};
\draw[arrow] (upper) -- (middle);
\draw[arrow] (middle) -- (lower);
\draw[arrow] (lower) -- (reviewer);
\draw[arrow] (reviewer.east) .. controls +(1.0,0.7) and +(1.0,-0.7) .. (upper.east);
\end{tikzpicture}
\caption{The four-agent stack mirrors the Quantum automation project, but the
contracts are SDE/Sampling objects rather than circuit oracles.}
\end{figure}

